\documentclass[12pt,a4paper]{article}

\usepackage[spanish,es-nodecimaldot]{babel}
\usepackage[utf8]{inputenc}
\usepackage[T1]{fontenc}
\usepackage{lmodern}
\usepackage{geometry}
\usepackage{microtype}
\usepackage{setspace}
\usepackage{parskip}
\usepackage{enumitem}
\usepackage{xcolor}
\usepackage{listings}
\usepackage{hyperref}
\usepackage{amsmath}
\usepackage{amssymb}
\usepackage{array}
\usepackage{booktabs}
\usepackage{longtable}

\geometry{margin=2.5cm}
\onehalfspacing
\setlist[itemize]{leftmargin=*,itemsep=0.2em}
\setlist[enumerate]{leftmargin=*,itemsep=0.2em}

\hypersetup{
  colorlinks=true,
  linkcolor=black,
  urlcolor=blue,
  pdftitle={Resolucion del Practico 4 - Testing de Software 2024},
  pdfauthor={}
}

\lstdefinestyle{codestyle}{
  basicstyle=\ttfamily\small,
  frame=single,
  breaklines=true,
  showstringspaces=false,
  keywordstyle=\color{blue!70!black}\bfseries,
  commentstyle=\color{gray!70!black}\itshape,
  stringstyle=\color{orange!70!black}
}

\title{\textbf{Resolucion del Practico 4}\\Testing de Software 2024}
\author{}
\date{}

\begin{document}

\maketitle
\tableofcontents
\newpage

\section*{Introduccion}
\addcontentsline{toc}{section}{Introduccion}

Este documento resuelve el \textit{Practico 4} de Testing de Software, centrado en el diseno de
tests mediante criterios estructurales sobre grafos: cobertura de nodos (NC), cobertura de aristas
(EC), cobertura de pares de aristas (EPC) y cobertura de caminos principales (PPC). Los dos
primeros ejercicios trabajan sobre grafos abstractos, mientras que los dos ultimos aplican los
mismos criterios sobre el Grafo de Flujo de Control (GCF) de los metodos \texttt{fmtRewrap} y
\texttt{patternIndex}.

A lo largo del documento se utilizan las siguientes convenciones:

\begin{itemize}
  \item \textbf{NC} (Node Coverage): cubrir cada nodo alcanzable.
  \item \textbf{EC} (Edge Coverage): cubrir cada camino de longitud $\le 1$ (cada arista).
  \item \textbf{EPC} (Edge-Pair Coverage): cubrir cada camino de longitud $\le 2$ (cada par de
  aristas consecutivas).
  \item \textbf{PPC} (Prime Path Coverage): cubrir todos los caminos simples maximales (los que
  no son subcamino propio de ningun otro camino simple).
\end{itemize}

\section{Ejercicio 1: Grafo de 4 nodos}

Se considera el grafo dirigido:
\begin{itemize}
  \item $N = \{1, 2, 3, 4\}$, $N_0 = \{1\}$, $N_f = \{4\}$.
  \item $E = \{(1,2), (2,3), (3,2), (2,4)\}$.
\end{itemize}

\textbf{Estructura.} Desde 1 solo se puede ir a 2; desde 2 se ramifica a 3 o a 4; desde 3 se
vuelve a 2; 4 es el unico nodo final y no tiene aristas salientes. Por consiguiente, todo camino
de test valido tiene la forma:
\[
1 \to 2 \to \underbrace{(\text{ciclo } 2 \leftrightarrow 3 \text{ repetido } k \ge 0 \text{ veces})}_{\text{opcional}} \to 4.
\]

Algunos caminos de test validos: $[1,2,4]$, $[1,2,3,2,4]$, $[1,2,3,2,3,2,4]$, etc.

Requisitos por criterio:
\begin{itemize}
  \item \textbf{NC}: $\{1, 2, 3, 4\}$.
  \item \textbf{EC}: $\{(1,2), (2,3), (3,2), (2,4)\}$.
  \item \textbf{EPC} (caminos de longitud 2 alcanzables): $[1,2,3]$, $[1,2,4]$, $[2,3,2]$,
  $[3,2,3]$, $[3,2,4]$.
\end{itemize}

\subsection{(a) NC pero no EC}

\textbf{No es posible.} La topologia del grafo fuerza que cualquier conjunto de tests que cubra
NC cubra automaticamente EC, por la siguiente razon:

\begin{enumerate}
  \item Visitar 1 obliga a tomar $(1,2)$ (unica arista saliente de 1).
  \item Visitar 3 obliga a tomar $(2,3)$ (unica arista entrante a 3).
  \item Una vez en 3, para poder llegar a un final hay que tomar $(3,2)$ (unica arista saliente).
  \item Para terminar en 4 hay que usar $(2,4)$ (unica arista entrante a 4).
\end{enumerate}

Los cuatro arcos asi forzados son exactamente $E$, por lo que NC implica EC en este grafo.

\subsection{(b) EC pero no EPC}

\textbf{Si es posible.} Tomamos $t_1 = [1, 2, 3, 2, 4]$.

\begin{itemize}
  \item Pares de aristas cubiertos por $t_1$: $[1,2,3]$, $[2,3,2]$, $[3,2,4]$.
  \item Aristas cubiertas: $(1,2)$, $(2,3)$, $(3,2)$, $(2,4)$, es decir, todo $E$.
\end{itemize}

Asi, $\{t_1\}$ satisface EC, pero deja sin cubrir $[1,2,4]$ y $[3,2,3]$, por lo que no satisface
EPC.

\subsection{(c) Conjunto que satisface EPC}

Un conjunto minimo es:
\begin{itemize}
  \item $t_1 = [1, 2, 4]$ \quad cubre $[1,2,4]$.
  \item $t_2 = [1, 2, 3, 2, 3, 2, 4]$ \quad cubre $[1,2,3]$, $[2,3,2]$, $[3,2,3]$, $[3,2,4]$.
\end{itemize}

La union cubre los cinco pares de aristas factibles. No alcanza con un unico test porque
$[1,2,3]$ y $[1,2,4]$ son mutuamente excluyentes en el prefijo: una vez que el test elige una de
las dos ramas en 2, no puede volver atras para tomar la otra.

\section{Ejercicio 2: Grafo de 7 nodos}

Grafo:
\begin{itemize}
  \item $N = \{1,2,3,4,5,6,7\}$, $N_0 = \{1\}$, $N_f = \{7\}$.
  \item $E = \{(1,2), (1,7), (2,3), (2,4), (3,2), (4,5), (4,6), (5,6), (6,1)\}$.
\end{itemize}

Caminos candidatos: $p_1 = [1,2,4,5,6,1,7]$, $p_2 = [1,2,3,2,4,6,1,7]$,
$p_3 = [1,2,3,2,4,5,6,1,7]$.

\subsection{(a) Requisitos de EPC (12)}

Los pares de aristas alcanzables son:

\begin{enumerate}
  \item $[1,2,3]$
  \item $[1,2,4]$
  \item $[2,3,2]$
  \item $[2,4,5]$
  \item $[2,4,6]$
  \item $[3,2,3]$
  \item $[3,2,4]$
  \item $[4,5,6]$
  \item $[4,6,1]$
  \item $[5,6,1]$
  \item $[6,1,2]$
  \item $[6,1,7]$
\end{enumerate}

\subsection{(b) ¿\{p1, p2, p3\} satisface EPC?}

\textbf{No.} Analizando los pares cubiertos por cada candidato:

\begin{itemize}
  \item $p_1$ cubre: $[1,2,4]$, $[2,4,5]$, $[4,5,6]$, $[5,6,1]$, $[6,1,7]$.
  \item $p_2$ cubre: $[1,2,3]$, $[2,3,2]$, $[3,2,4]$, $[2,4,6]$, $[4,6,1]$, $[6,1,7]$.
  \item $p_3$ cubre: $[1,2,3]$, $[2,3,2]$, $[3,2,4]$, $[2,4,5]$, $[4,5,6]$, $[5,6,1]$, $[6,1,7]$.
\end{itemize}

Union: $\{[1,2,3], [1,2,4], [2,3,2], [2,4,5], [2,4,6], [3,2,4], [4,5,6], [4,6,1], [5,6,1],
[6,1,7]\}$.

\textbf{Pares faltantes:} $[3,2,3]$ y $[6,1,2]$. Ninguno de los tres caminos repite el ciclo
$2 \to 3 \to 2$ ni retoma el ciclo externo despues de volver a 1.

\subsection{(c) Tour directo o con sidetrips}

Camino simple $q = [3, 2, 4, 5, 6]$; camino de test $t = [1, 2, 3, 2, 4, 6, 1, 2, 4, 5, 6, 1, 7]$.

\textbf{No es tour directo:} la primera vez que $t$ llega a 4 toma $4 \to 6$, mientras que $q$
exige $4 \to 5$. Por lo tanto las aristas de $q$ no aparecen \emph{en el mismo orden} sin
interrupciones.

\textbf{Si es tour con sidetrip.} La idea de un sidetrip es: desde un nodo del camino simple, el
camino de test se desvia y vuelve al \emph{mismo} nodo del camino simple. Concretamente:

\begin{itemize}
  \item $t$ pasa por 3, luego por 2, luego por 4 (tres primeros nodos de $q$).
  \item Desde 4, $t$ se desvia tomando $[4, 6, 1, 2, 4]$ y vuelve al mismo nodo 4.
  \item Recien entonces continua con $4 \to 5 \to 6$, completando los dos nodos finales de $q$.
\end{itemize}

\textbf{Sidetrip:} $[4, 6, 1, 2, 4]$.

\subsection{(d) Requisitos para NC, EC y PPC}

\textbf{NC:} $\{1, 2, 3, 4, 5, 6, 7\}$.

\textbf{EC:} las nueve aristas $(1,2), (1,7), (2,3), (2,4), (3,2), (4,5), (4,6), (5,6), (6,1)$.

\textbf{PPC (caminos principales).} Un camino principal es un camino simple maximal. Los caminos
simples relevantes (no extensibles a otro camino simple) son:

\begin{enumerate}
  \item $[1, 2, 4, 6, 1]$
  \item $[1, 2, 4, 5, 6, 1]$
  \item $[2, 3, 2]$
  \item $[2, 4, 6, 1, 2]$
  \item $[2, 4, 5, 6, 1, 2]$
  \item $[3, 2, 3]$
  \item $[3, 2, 4, 6, 1, 7]$
  \item $[3, 2, 4, 5, 6, 1, 7]$
  \item $[4, 6, 1, 2, 3]$
  \item $[4, 6, 1, 2, 4]$
  \item $[4, 5, 6, 1, 2, 3]$
  \item $[4, 5, 6, 1, 2, 4]$
  \item $[5, 6, 1, 2, 4, 5]$
  \item $[6, 1, 2, 4, 6]$
  \item $[6, 1, 2, 4, 5, 6]$
\end{enumerate}

\subsection{(e) NC pero no EC}

Posible con $t = [1, 2, 3, 2, 4, 5, 6, 1, 7]$: visita los siete nodos, pero no recorre $(4, 6)$
(al llegar a 4 toma la rama hacia 5). Por lo tanto cumple NC pero falla EC.

\subsection{(f) EC pero no PPC}

Posible con $\{t_1, t_2\}$ donde $t_1 = [1, 2, 4, 5, 6, 1, 7]$ y $t_2 = [1, 2, 3, 2, 4, 6, 1, 7]$.

\begin{itemize}
  \item Aristas cubiertas: $(1,2), (2,4), (4,5), (5,6), (6,1), (1,7)$ (de $t_1$) y $(2,3), (3,2),
  (4,6)$ (de $t_2$). Las nueve aristas estan cubiertas.
  \item Pero el prime path $[3, 2, 3]$ no aparece: ningun test repite el ciclo
  $2 \to 3 \to 2 \to 3$. Tampoco aparece $[6, 1, 2, 4, 6]$ (no se ejecuta el ciclo externo dos
  veces consecutivas), entre otros.
\end{itemize}

Como queda al menos un prime path sin cubrir, el conjunto satisface EC pero no PPC.

\section{Ejercicio 3: \texttt{fmtRewrap}}

Se trabaja sobre el metodo \texttt{fmtRewrap(String S, int N)} de
\texttt{assignment4\_exercises.FmtRewrap.java}. Para mantener el analisis manejable se utiliza un
CFG \emph{por bloques} agrupando las sentencias logicas del cuerpo del \texttt{while}.

\subsection{Modelo de CFG (a)}

\textbf{Nodos del modelo.}
\begin{itemize}
  \item \textbf{W}: condicion del \texttt{while} (\texttt{i < S.length()}).
  \item \textbf{C}: clasificacion del caracter actual (cadena \texttt{if/else if} que asigna
  \texttt{state}).
  \item \textbf{BW}: caso \texttt{betweenWord} dentro del \texttt{switch}.
  \item \textbf{LB}: caso \texttt{lineBreak} dentro del \texttt{switch}.
  \item \textbf{CR?}: caso \texttt{crFound} (decide entre hard-CR y soft-CR).
  \item \textbf{CRh}: rama hard-CR (\texttt{i+1 < S.length() \&\& SArr[i+1] == CR}).
  \item \textbf{CRs}: rama soft-CR (rama \texttt{else} del caso anterior).
  \item \textbf{IW}: caso \texttt{inWord} o por defecto.
  \item \textbf{INC}: incremento \texttt{i++}.
  \item \textbf{EXIT}: salida del \texttt{while} y \texttt{return}.
\end{itemize}

\textbf{Aristas:}
\[
\begin{aligned}
&\mathbf{W} \to \mathbf{C} \quad (\text{true}), \qquad \mathbf{W} \to \mathbf{EXIT} \quad (\text{false}), \\
&\mathbf{C} \to \mathbf{BW},\quad \mathbf{C} \to \mathbf{LB},\quad \mathbf{C} \to \mathbf{CR?},\quad \mathbf{C} \to \mathbf{IW}, \\
&\mathbf{BW} \to \mathbf{INC},\quad \mathbf{LB} \to \mathbf{INC}, \\
&\mathbf{CR?} \to \mathbf{CRh},\quad \mathbf{CR?} \to \mathbf{CRs}, \\
&\mathbf{CRh} \to \mathbf{INC},\quad \mathbf{CRs} \to \mathbf{INC},\quad \mathbf{IW} \to \mathbf{INC}, \\
&\mathbf{INC} \to \mathbf{W}.
\end{aligned}
\]

\subsection{(b) Test que sale del \texttt{while} sin entrar al cuerpo}

\textbf{Caso:} $t = (\texttt{S}=\texttt{""}, N=10)$.

Cuando \texttt{S} es la cadena vacia, \texttt{S.length() == 0} y la guarda del \texttt{while}
\texttt{i < S.length()} es \texttt{0 < 0 = false}. El flujo toma directamente
$\mathbf{W} \to \mathbf{EXIT}$ y ejecuta \texttt{S = new String(SArr) + CR; return}, devolviendo
\texttt{"\textbackslash n"}.

\subsection{(c) Requisitos NC, EC y PPC}

\textbf{NC:} $\{\mathbf{W}, \mathbf{C}, \mathbf{BW}, \mathbf{LB}, \mathbf{CR?}, \mathbf{CRh},
\mathbf{CRs}, \mathbf{IW}, \mathbf{INC}, \mathbf{EXIT}\}$ (10 requisitos).

\textbf{EC:} las 14 aristas listadas en el CFG.

\textbf{PPC:} sobre este CFG reducido se obtienen 47 prime paths, agrupados por su nodo inicial:

\begin{enumerate}
  \item $[W, C, BW, INC, W]$
  \item $[W, C, LB, INC, W]$
  \item $[W, C, IW, INC, W]$
  \item $[W, C, CR?, CRh, INC, W]$
  \item $[W, C, CR?, CRs, INC, W]$
  \item $[C, BW, INC, W, C]$
  \item $[C, BW, INC, W, EXIT]$
  \item $[C, LB, INC, W, C]$
  \item $[C, LB, INC, W, EXIT]$
  \item $[C, IW, INC, W, C]$
  \item $[C, IW, INC, W, EXIT]$
  \item $[C, CR?, CRh, INC, W, C]$
  \item $[C, CR?, CRh, INC, W, EXIT]$
  \item $[C, CR?, CRs, INC, W, C]$
  \item $[C, CR?, CRs, INC, W, EXIT]$
  \item $[BW, INC, W, C, BW]$
  \item $[BW, INC, W, C, LB]$
  \item $[BW, INC, W, C, IW]$
  \item $[BW, INC, W, C, CR?, CRh]$
  \item $[BW, INC, W, C, CR?, CRs]$
  \item $[LB, INC, W, C, BW]$
  \item $[LB, INC, W, C, LB]$
  \item $[LB, INC, W, C, IW]$
  \item $[LB, INC, W, C, CR?, CRh]$
  \item $[LB, INC, W, C, CR?, CRs]$
  \item $[CR?, CRh, INC, W, C, BW]$
  \item $[CR?, CRh, INC, W, C, LB]$
  \item $[CR?, CRh, INC, W, C, CR?]$
  \item $[CR?, CRh, INC, W, C, IW]$
  \item $[CR?, CRs, INC, W, C, BW]$
  \item $[CR?, CRs, INC, W, C, LB]$
  \item $[CR?, CRs, INC, W, C, CR?]$
  \item $[CR?, CRs, INC, W, C, IW]$
  \item $[CRh, INC, W, C, CR?, CRh]$
  \item $[CRh, INC, W, C, CR?, CRs]$
  \item $[CRs, INC, W, C, CR?, CRh]$
  \item $[CRs, INC, W, C, CR?, CRs]$
  \item $[IW, INC, W, C, BW]$
  \item $[IW, INC, W, C, LB]$
  \item $[IW, INC, W, C, IW]$
  \item $[IW, INC, W, C, CR?, CRh]$
  \item $[IW, INC, W, C, CR?, CRs]$
  \item $[INC, W, C, BW, INC]$
  \item $[INC, W, C, LB, INC]$
  \item $[INC, W, C, IW, INC]$
  \item $[INC, W, C, CR?, CRh, INC]$
  \item $[INC, W, C, CR?, CRs, INC]$
\end{enumerate}

\subsection{(d) Suite para NC (pero no EC)}

Suite: \texttt{FmtRewrapNodeCoverageTest.java}, ejecutable con
\begin{lstlisting}[style=codestyle]
JAVA_HOME=$(/usr/libexec/java_home -v 17) mvn -Dmaven.repo.local=.m2 \
    -Dtest=FmtRewrapNodeCoverageTest test
\end{lstlisting}

JaCoCo reporta sobre \texttt{fmtRewrap}: \texttt{LINE 29/29}, \texttt{BRANCH 16/16}.

\textbf{Observacion.} Para esta implementacion y este nivel de modelado, cubrir todas las
sentencias termina cubriendo tambien todas las ramas que la herramienta mide. En la practica,
\emph{no} se logro armar una suite que satisfaga NC sin satisfacer EC al mismo tiempo segun
JaCoCo. La razon es que muchas sentencias y ramas estan tan acopladas que llegar a un nodo del
CFG fuerza la ejecucion de su arista saliente unica.

\subsection{(e) Suite para EC pero no PPC}

Suite: \texttt{FmtRewrapEdgeButNotPrimePathCoverageTest.java}. JaCoCo reporta nuevamente
\texttt{LINE 29/29} y \texttt{BRANCH 16/16}.

\textbf{Por que la suite no satisface PPC:} no se recorre, por ejemplo, el prime path
$[BW, INC, W, C, IW]$. Intuitivamente, cada vez que se entra en \textbf{BW} (caso
\texttt{betweenWord}), la siguiente iteracion ``util'' del bucle pasa por \textbf{LB} o por otra
rama, no por \textbf{IW}, por lo que esa secuencia particular nunca aparece en los caminos
ejecutados.

\subsection{(f) Suite para PPC con Best Effort Touring}

Suite: \texttt{FmtRewrapPrimePathBestEffortCoverageTest.java}. JaCoCo: \texttt{LINE 29/29},
\texttt{BRANCH 16/16}.

\textbf{Cobertura sobre el CFG modelado.}
\begin{itemize}
  \item Prime paths totales: 47.
  \item Cubiertos directamente: 43.
  \item No cubiertos: 4. Todos \emph{infactibles} por la semantica del metodo:
  \begin{enumerate}
    \item $[CR?, CRh, INC, W, C, LB]$
    \item $[CR?, CRs, INC, W, C, CR?]$
    \item $[CRs, INC, W, C, CR?, CRh]$
    \item $[CRs, INC, W, C, CR?, CRs]$
  \end{enumerate}
\end{itemize}

\textbf{Justificacion de la infactibilidad:}
\begin{itemize}
  \item \textbf{Hard-CR seguido de LB.} Despues de tomar CRh, el codigo fija \texttt{col = 1}.
  Para que en la siguiente iteracion se entre a LB seria necesario que \texttt{col >= N}, pero
  \texttt{N} debiera ser tan pequeno que la iteracion previa ya habria disparado LB y no se
  habria podido tomar CRh. Las dos condiciones son mutuamente excluyentes.

  \item \textbf{Soft-CR seguido de otro CR?.} CRs se elige justamente cuando el caracter
  siguiente \emph{no} es CR. En la iteracion siguiente, el caracter actual ya no es CR, por lo
  que no se puede volver a entrar a CR? en la posicion inmediata.
\end{itemize}

\textbf{Conclusion.} La suite cubre todos los prime paths factibles del CFG bajo \textit{Best
Effort Touring} y mantiene cobertura total de lineas y ramas. PPC subsume EC, por lo que la
cobertura de aristas tambien queda garantizada.

\section{Ejercicio 4: \texttt{patternIndex} -- instrumentacion y cobertura}

Se trabaja sobre \texttt{patternIndex(String subject, String pattern)} en
\texttt{assignment4\_exercises.PatternIndex.java}.

\subsection{Instrumentacion}

Para registrar los caminos efectivamente ejecutados, se inserto una llamada de tipo
\texttt{PatternIndexPathTracker.hit("X")} en cada nodo del CFG. La clase auxiliar
\texttt{PatternIndexPathTracker} esta en
\texttt{src/main/java/assignment4\_exercises/PatternIndexPathTracker.java} y mantiene, por
invocacion, la secuencia de nodos visitados.

\textbf{Nodos instrumentados:}
\begin{itemize}
  \item \textbf{S}: inicio del metodo.
  \item \textbf{W}: condicion del \texttt{while} externo.
  \item \textbf{I}: comparacion del primer caracter del patron.
  \item \textbf{M}: bloque de inicializacion cuando coincide el primer caracter.
  \item \textbf{FT}: condicion verdadera del \texttt{for} interno.
  \item \textbf{IFM}: chequeo de mismatch dentro del \texttt{for}.
  \item \textbf{MIS}: bloque de mismatch seguido de \texttt{break}.
  \item \textbf{FI}: avance del \texttt{for} cuando no hubo mismatch.
  \item \textbf{FF}: terminacion del \texttt{for} sin \texttt{break}.
  \item \textbf{INC}: incremento \texttt{iSub++}.
  \item \textbf{R}: \texttt{return}.
\end{itemize}

\textbf{Aristas del CFG instrumentado:}
\[
\begin{aligned}
&S \to W,\quad W \to I,\quad W \to R, \\
&I \to M,\quad I \to INC, \\
&M \to FT,\quad M \to FF, \\
&FT \to IFM,\quad IFM \to MIS,\quad IFM \to FI, \\
&FI \to FT,\quad FI \to FF, \\
&MIS \to INC,\quad FF \to INC,\quad INC \to W.
\end{aligned}
\]
En total 15 aristas.

\subsection{Suite JUnit de la tabla del enunciado}

La suite \texttt{PatternIndexInstrumentationCoverageTest.java} ejecuta los 10 casos pedidos:

\begin{longtable}{@{}p{2cm}p{2cm}p{1.5cm}@{}}
\toprule
\textbf{subject} & \textbf{pattern} & \textbf{output} \\
\midrule
\endhead
\texttt{"a"} & \texttt{"bc"} & $-1$ \\
\texttt{"ab"} & \texttt{"a"} & $0$ \\
\texttt{"ab"} & \texttt{"ab"} & $0$ \\
\texttt{"ab"} & \texttt{"ac"} & $-1$ \\
\texttt{"ab"} & \texttt{"b"} & $1$ \\
\texttt{"ab"} & \texttt{"c"} & $-1$ \\
\texttt{"abc"} & \texttt{"abc"} & $0$ \\
\texttt{"abc"} & \texttt{"abd"} & $-1$ \\
\texttt{"abc"} & \texttt{"ba"} & $-1$ \\
\texttt{"abc"} & \texttt{"bc"} & $1$ \\
\bottomrule
\end{longtable}

Al finalizar la suite, un metodo \texttt{@AfterClass} consulta el tracker, calcula automaticamente
los caminos principales del CFG por DFS y genera un reporte con:

\begin{itemize}
  \item secuencia de nodos por cada invocacion,
  \item cobertura de aristas (cubiertas / total),
  \item cobertura de caminos principales (cubiertos / total),
  \item lista de requisitos cubiertos y faltantes por criterio.
\end{itemize}

Ubicacion del reporte: \texttt{tp4/assignmnet-4-rodeghiero/target/patternindex-path-report.txt}.

\subsection{Ejecucion}

\begin{lstlisting}[style=codestyle]
JAVA_HOME=$(/usr/libexec/java_home -v 17) mvn -Dmaven.repo.local=.m2 \
    -Dtest=PatternIndexInstrumentationCoverageTest test
\end{lstlisting}

\subsection{Resultados}

Sobre el CFG instrumentado y con los 10 casos de la tabla:

\begin{itemize}
  \item \textbf{Cobertura de aristas:} $15 / 15$ ($100\%$).
  \item \textbf{Cobertura de caminos principales:} $22 / 39$.
\end{itemize}

\textbf{Interpretacion.}
\begin{itemize}
  \item Los 10 casos son \emph{suficientes para saturar EC}: cada decision del metodo termina
  ejercitandose por lo menos una vez (match al inicio, mismatch al medio, mismatch al final,
  patron mas largo que el sujeto, etc.).
  \item En cambio, no se alcanza PPC. Esto era esperable: PPC subsume EC pero no al reves. Los
  prime paths faltantes corresponden a combinaciones de iteraciones del \texttt{while} externo y
  del \texttt{for} interno que esos diez casos puntuales no llegan a ejercitar (por ejemplo,
  patrones que matchean parcialmente en una posicion intermedia, fallan y luego vuelven a
  matchear mas adelante).
  \item La instrumentacion permite ver exactamente, por cada caso, que ruta del CFG se ejecuto
  y por que un prime path queda fuera: el reporte enumera todos los requisitos faltantes con su
  nombre canonico (lista de nodos).
\end{itemize}

\section*{Conclusiones}
\addcontentsline{toc}{section}{Conclusiones}

El practico expone con claridad la jerarquia de subsuncion
$\text{PPC} \supseteq \text{EPC} \supseteq \text{EC} \supseteq \text{NC}$ y muestra, mediante
ejemplos concretos, por que cada nivel ``extra'' aporta requisitos distintos:

\begin{itemize}
  \item En el grafo del Ejercicio 1, NC e EC \emph{coinciden} en la practica porque la topologia
  fuerza cada nodo a usar todas sus aristas. EC y EPC se distinguen recien con un test minimo.

  \item En el grafo del Ejercicio 2 aparece la nocion de \textit{sidetrip}: un camino de test
  puede ``visitar'' un camino simple aunque no lo recorra de manera contigua, siempre que se
  vuelva al mismo nodo del que se desvio.

  \item Sobre codigo real (\texttt{fmtRewrap} y \texttt{patternIndex}), los criterios estructurales
  funcionan como una guia para disenar tests con \emph{intencion} y para medir, con herramientas
  como JaCoCo o con instrumentacion propia, cuanto del comportamiento del programa queda
  efectivamente ejercitado. \textit{Best Effort Touring} resulta especialmente util cuando la
  semantica del programa hace infactibles algunos prime paths.
\end{itemize}

\end{document}
